% GENERATED FILE. Edit canonical lessons, then run npm run generate:books.
% canonical-source-hash: fnv1a64:0e7f2bacba5ac826
% canonical-lessons: RU-C529-pyl, RU-C529-uborka, RU-C529-pylesosit, RU-C529-tryapka, RU-C529-vytirat

\chapter{Cleaning Day}
\label{ch:ru-cleaning-day}
\hlchaptermodality{529}

\begin{chapteropening}
\textbf{By the end of this chapter:} \emph{I can talk about cleaning the house.}

Complete the last lesson of chapter 529: I can talk about cleaning the house.
\end{chapteropening}

\section[\texorpdfstring{\ru{пыль} (pyl)}{пыль (pyl)}]{\ru{пыль} (pyl) — dust}

\label{lesson:RU-C529-pyl}



\begin{quote}
Before the new one: say the Russian for a subscriber, then the Russian for a gamer.
\end{quote}

\subsection*{You'll want to know: \ru{пыль}}
\textbf{\ru{пыль}} (\emph{pyl}) — \textquotedblleft{}dust\textquotedblright{}.

\begin{quote}
\ru{На} \ru{полках} \ru{много} \ru{пыли}.
\end{quote}

\begin{quote}
\emph{Na pólkakh mnógo pýli.}
\end{quote}

\begin{quote}
There's a lot of dust on the shelves.
\end{quote}

Feminine; no plural.

\begin{cousinweb}
\emph{\ru{пыль}} is inherited from Proto-Slavic \emph{pyl\ru{ъ}} (EtymDB).
\end{cousinweb}

\begin{grammarlens}[title={What you've built}]
One more word for telling, explaining and describing at B1.
\end{grammarlens}

\subsection*{Your turn}
\begin{itemize}
\raggedright
  \item \emph{Say it:} \emph{pyl}
  \item \emph{Say it:} \emph{pyl}, once more
  \item \emph{Say it:} a sentence of your own with \emph{pyl}
  \item \emph{Recall:} say the Russian for a subscriber, then the Russian for a gamer, then say \emph{pyl} again
\end{itemize}

\begin{tcolorbox}[breakable,colback=teal!4,colframe=teal!35!black,title={Before you move on}]
What does \ru{пыль} mean? (\textquotedblleft{}Dust\textquotedblright{}.) Say it once more.
\end{tcolorbox}
\section[\texorpdfstring{\ru{уборка} (ubórka)}{уборка (ubórka)}]{\ru{уборка} (ubórka) — cleaning, tidying up}

\label{lesson:RU-C529-uborka}



\begin{quote}
Before the new one: say the Russian for a gamer, then the Russian for dust.
\end{quote}

\subsection*{You'll want to know: \ru{уборка}}
\textbf{\ru{уборка}} (\emph{ubórka}) — \textquotedblleft{}cleaning, tidying up\textquotedblright{}.

\begin{quote}
\ru{В} \ru{субботу} \ru{у} \ru{нас} \ru{генеральная} \ru{уборка}.
\end{quote}

\begin{quote}
\emph{V subbótu u nas generálnaya ubórka.}
\end{quote}

\begin{quote}
On Saturday we're having a big clean.
\end{quote}

Feminine. \emph{\ru{Генеральная} \ru{уборка}} — a thorough clean.

\begin{cousinweb}
\emph{\ru{уборка}} is built on \emph{\ru{убрать}} 'to tidy away', from \emph{\ru{у}-} and \emph{\ru{брать}} 'to take'; EtymDB lists \emph{\ru{уборка}} without a source.
\end{cousinweb}

\begin{grammarlens}[title={What you've built}]
One more word for telling, explaining and describing at B1.
\end{grammarlens}

\subsection*{Your turn}
\begin{itemize}
\raggedright
  \item \emph{Say it:} \emph{ubórka}
  \item \emph{Say it:} \emph{ubórka}, once more
  \item \emph{Say it:} a sentence of your own with \emph{ubórka}
  \item \emph{Recall:} say the Russian for a gamer, then the Russian for dust, then say \emph{ubórka} again
\end{itemize}

\begin{tcolorbox}[breakable,colback=teal!4,colframe=teal!35!black,title={Before you move on}]
What does \ru{уборка} mean? (\textquotedblleft{}Cleaning, tidying up\textquotedblright{}.) Say it once more.
\end{tcolorbox}
\section[\texorpdfstring{\ru{пылесосить} (pylesósit)}{пылесосить (pylesósit)}]{\ru{пылесосить} (pylesósit) — to vacuum}

\label{lesson:RU-C529-pylesosit}



\begin{quote}
Before the new one: say the Russian for dust, then the Russian for cleaning.
\end{quote}

\subsection*{You'll want to know: \ru{пылесосить}}
\textbf{\ru{пылесосить}} (\emph{pylesósit}) — \textquotedblleft{}to vacuum\textquotedblright{}.

\begin{quote}
\ru{Сын} \ru{пылесосит} \ru{ковёр} \ru{в} \ru{гостиной}.
\end{quote}

\begin{quote}
\emph{Syn pylesósit kovyór v gostínoy.}
\end{quote}

\begin{quote}
My son is vacuuming the carpet in the living room.
\end{quote}

Imperfective; present \ru{пылесошу}, \ru{пылесосишь} (\ru{с} becomes \ru{ш}). Perfective \ru{пропылесосить}.

\begin{cousinweb}
\emph{\ru{пылесосить}} is built on \emph{\ru{пылесос}} 'a vacuum cleaner', a compound of \emph{\ru{пыль}} 'dust' and \emph{\ru{сосать}} 'to suck'; EtymDB has no entry.
\end{cousinweb}

\begin{grammarlens}[title={What you've built}]
One more word for telling, explaining and describing at B1.
\end{grammarlens}

\subsection*{Your turn}
\begin{itemize}
\raggedright
  \item \emph{Say it:} \emph{pylesósit}
  \item \emph{Say it:} \emph{pylesósit}, once more
  \item \emph{Say it:} a sentence of your own with \emph{pylesósit}
  \item \emph{Recall:} say the Russian for dust, then the Russian for cleaning, then say \emph{pylesósit} again
\end{itemize}

\begin{tcolorbox}[breakable,colback=teal!4,colframe=teal!35!black,title={Before you move on}]
What does \ru{пылесосить} mean? (\textquotedblleft{}To vacuum\textquotedblright{}.) Say it once more.
\end{tcolorbox}
\section[\texorpdfstring{\ru{тряпка} (tryápka)}{тряпка (tryápka)}]{\ru{тряпка} (tryápka) — a cloth, a rag}

\label{lesson:RU-C529-tryapka}



\begin{quote}
Before the new one: say the Russian for cleaning, then the Russian for to vacuum.
\end{quote}

\subsection*{You'll want to know: \ru{тряпка}}
\textbf{\ru{тряпка}} (\emph{tryápka}) — \textquotedblleft{}a cloth, a rag\textquotedblright{}.

\begin{quote}
\ru{Мама} \ru{вытерла} \ru{кухонный} \ru{стол} \ru{мокрой} \ru{тряпкой}.
\end{quote}

\begin{quote}
\emph{Máma výterla kúkhonnyy stol mókroy tryápkoy.}
\end{quote}

\begin{quote}
Mum wiped the kitchen table with a damp cloth.
\end{quote}

Feminine.

\begin{cousinweb}
EtymDB lists \emph{\ru{тряпка}} without a source; it is usually linked to \emph{\ru{тряпьё}} 'rags', and the further origin is not settled.
\end{cousinweb}

\begin{grammarlens}[title={What you've built}]
One more word for telling, explaining and describing at B1.
\end{grammarlens}

\subsection*{Your turn}
\begin{itemize}
\raggedright
  \item \emph{Say it:} \emph{tryápka}
  \item \emph{Say it:} \emph{tryápka}, once more
  \item \emph{Say it:} a sentence of your own with \emph{tryápka}
  \item \emph{Recall:} say the Russian for cleaning, then the Russian for to vacuum, then say \emph{tryápka} again
\end{itemize}

\begin{tcolorbox}[breakable,colback=teal!4,colframe=teal!35!black,title={Before you move on}]
What does \ru{тряпка} mean? (\textquotedblleft{}A cloth, a rag\textquotedblright{}.) Say it once more.
\end{tcolorbox}
\section[\texorpdfstring{\ru{вытирать} (vytirát)}{вытирать (vytirát)}]{\ru{вытирать} (vytirát) — to wipe; to dry}

\label{lesson:RU-C529-vytirat}



\begin{quote}
Before the new one: say the Russian for to vacuum, then the Russian for a cloth.
\end{quote}

\subsection*{You'll want to know: \ru{вытирать}}
\textbf{\ru{вытирать}} (\emph{vytirát}) — \textquotedblleft{}to wipe; to dry\textquotedblright{}.

\begin{quote}
\ru{После} \ru{ужина} \ru{брат} \ru{моет} \ru{тарелки}, \ru{а} \ru{я} \ru{их} \ru{вытираю}.
\end{quote}

\begin{quote}
\emph{Pósle úzhina brat móyet tarélki, a ya ikh vytiráyu.}
\end{quote}

\begin{quote}
After dinner my brother washes the plates and I dry them.
\end{quote}

Imperfective; perfective \ru{вытереть} (\ru{вытру}; past \ru{вытер}). \emph{\ru{Вытирать} \ru{пыль}} — to dust.

\begin{cousinweb}
\emph{\ru{вытирать}} is built on \emph{\ru{вы}-} and the root of \emph{\ru{тереть}} 'to rub'; EtymDB has no entry.
\end{cousinweb}

\begin{grammarlens}[title={What you've built}]
One more word for telling, explaining and describing at B1.
\end{grammarlens}

\subsection*{Your turn}
\begin{itemize}
\raggedright
  \item \emph{Say it:} \emph{vytirát}
  \item \emph{Say it:} \emph{vytirát}, once more
  \item \emph{Say it:} a sentence of your own with \emph{vytirát}
  \item \emph{Recall:} say the Russian for to vacuum, then the Russian for a cloth, then say \emph{vytirát} again
\end{itemize}

\begin{tcolorbox}[breakable,colback=teal!4,colframe=teal!35!black,title={Before you move on}]
What does \ru{вытирать} mean? (\textquotedblleft{}To wipe; to dry\textquotedblright{}.) Say it once more.
\end{tcolorbox}
